%=================================================
% CHAPITRE 7
%=================================================

\chapter{Sprint 5 : Intégration, qualité et livraison}

\section{Introduction}

Le cinquième sprint correspond à la phase de consolidation du projet. Après avoir développé le socle d'authentification, la gouvernance d'accès, l'édition documentaire et les modules métier, il devient nécessaire de vérifier la cohérence globale de l'application. Cette étape vise à transformer un ensemble de fonctionnalités développées progressivement en une solution stable, testable et présentable dans le cadre d'un projet de fin d'études.

L'intégration finale couvre plusieurs dimensions : intégration frontend/backend, validation des parcours utilisateurs, tests fonctionnels, tests d'intégration, optimisation des performances, renforcement de la sécurité et préparation de la documentation. Elle ne consiste donc pas seulement à corriger des anomalies. Elle permet également d'améliorer l'expérience utilisateur, de fiabiliser les échanges API et de préparer la démonstration finale.

Ce sprint a également une importance académique. Il permet de montrer que le projet ne se limite pas à une implémentation technique, mais suit une démarche de qualité logicielle. Les tests, la documentation, les captures d'écran, les diagrammes et la préparation de la livraison contribuent à rendre le travail évaluable, reproductible et maintenable.

\section{Objectifs du sprint}

Les objectifs du sprint 5 sont les suivants :

\begin{itemize}
    \item intégrer et valider les flux complets entre Angular et ASP.NET Core ;
    \item vérifier les parcours principaux des trois rôles : super administrateur, administrateur et utilisateur final ;
    \item tester les opérations CRUD sur les familles, modèles, chartes graphiques et configurations de vues de tables ;
    \item valider le processus de génération documentaire dynamique ;
    \item améliorer la gestion des erreurs et l'affichage des messages utilisateur ;
    \item renforcer les contrôles de sécurité sur les endpoints sensibles ;
    \item optimiser le chargement des ressources et la réactivité des interfaces ;
    \item préparer la documentation technique, la documentation utilisateur et la démonstration PFE.
\end{itemize}

\section{Backlog du sprint}

\begin{longtable}{p{0.17\textwidth} p{0.38\textwidth} p{0.35\textwidth}}
\caption{Backlog du sprint 5}\\
\toprule
\textbf{User story} & \textbf{Description} & \textbf{Tâches principales}\\
\midrule
\endfirsthead
\toprule
\textbf{User story} & \textbf{Description} & \textbf{Tâches principales}\\
\midrule
\endhead
US5.1 & En tant qu'utilisateur final, je peux utiliser l'application sans blocage majeur. & Tester connexion, navigation, sélection famille, sélection modèle, prévisualisation et génération.\\
\midrule
US5.2 & En tant qu'administrateur, je peux gérer les ressources métier de bout en bout. & Tester CRUD modèles, chartes, familles et vues de tables.\\
\midrule
US5.3 & En tant que système, je dois afficher des erreurs compréhensibles. & Améliorer l'interceptor d'erreur, valider les réponses API, afficher les messages utiles.\\
\midrule
US5.4 & En tant qu'évaluateur PFE, je peux comprendre et reproduire l'exécution. & Rédiger documentation, préparer captures, produire diagrammes et décrire l'architecture.\\
\midrule
US5.5 & En tant que développeur, je peux livrer une version stable. & Exécuter build, tests prioritaires, revue des endpoints, vérification sécurité.\\
\bottomrule
\end{longtable}

\subsection{Intégration frontend/backend}

L'intégration frontend/backend consiste à vérifier que les contrats HTTP sont cohérents. Les services Angular doivent appeler les bons endpoints, envoyer les données sous la forme attendue et interpréter correctement les réponses. Côté backend, les contrôleurs doivent retourner des statuts HTTP cohérents et des messages d'erreur exploitables.

\subsection{Tests fonctionnels et tests d'intégration}

Les tests fonctionnels couvrent les scénarios visibles par l'utilisateur. Les tests d'intégration vérifient le fonctionnement conjoint des composants : formulaire Angular, service Angular, API, service backend, repository et base SQL Server. Cette distinction permet de détecter aussi bien les erreurs d'interface que les erreurs de contrat ou de persistance.

\subsection{Optimisation des performances}

Les optimisations portent principalement sur le chargement Angular, la limitation des données retournées, la pagination ou limitation des listes et la réduction des calculs inutiles lors du rendu documentaire. Les vues de tables utilisent par exemple une limite de lignes afin d'éviter de charger des volumes trop importants.

\subsection{Préparation de la livraison}

La préparation de la livraison inclut la documentation, les captures d'écran, les diagrammes, la vérification des prérequis techniques et la préparation d'un scénario de démonstration. L'objectif est de présenter une application cohérente et compréhensible par un jury.

\section{Analyse et conception}

\subsection{Diagramme de cas d'utilisation}

\diagramplaceholder{sprint5_cas_utilisation}{Diagramme de cas d'utilisation du sprint 5}{Cas d'utilisation d'intégration, qualité et livraison}

Le diagramme de cas d'utilisation du sprint 5 présente les activités transversales : valider un parcours complet, exécuter les tests, vérifier la sécurité, préparer la documentation et réaliser la démonstration. Les acteurs sont l'utilisateur final, l'administrateur, le super administrateur, le développeur et l'évaluateur PFE.

Ce diagramme montre que la livraison ne concerne pas uniquement le développeur. Elle implique aussi les utilisateurs qui valident les parcours, l'administrateur qui contrôle les données métier et l'évaluateur qui doit comprendre l'architecture et reproduire les scénarios.

\subsection{Diagramme de classes}

\diagramplaceholder{sprint5_classes}{Diagramme de classes du sprint 5}{Classes et services transversaux de la phase d'intégration}

Le diagramme de classes du sprint 5 met en évidence les composants transversaux : \texttt{ApiService}, \texttt{AuthInterceptor}, \texttt{ErrorInterceptor}, \texttt{EditorStateService}, \texttt{TemplateService}, \texttt{FamilyService}, \texttt{TableViewService} et \texttt{DocumentRenderService}. Ces services constituent la colonne vertébrale de l'application Angular.

Côté backend, \texttt{EditorController}, \texttt{AuthController}, \texttt{EditorService} et \texttt{EditorRepository} assurent l'orchestration des fonctionnalités. La phase d'intégration vérifie que ces composants collaborent correctement sur les parcours complets.

\subsection{Diagrammes de séquence}

\diagramplaceholder{sprint5_sequences}{Diagramme de séquence du sprint 5}{Séquence complète de démonstration finale}

Le diagramme de séquence décrit un parcours de démonstration complet. L'utilisateur se connecte, le frontend récupère l'état applicatif, l'administrateur vérifie les modèles et chartes, l'utilisateur sélectionne une famille, choisit un bénéficiaire, génère une prévisualisation, puis consulte le résultat. Chaque étape traverse plusieurs couches et constitue donc un test d'intégration réel.

\subsubsection*{Description des interactions système}

Le sprint 5 vérifie la continuité entre les modules développés. Un problème dans l'authentification empêcherait le chargement de l'état. Une erreur dans les familles empêcherait la sélection du bénéficiaire. Une incohérence dans les modèles provoquerait un rendu incomplet. La phase d'intégration consiste donc à tester les dépendances fonctionnelles dans leur ordre réel d'utilisation.

\subsubsection*{Analyse technique}

L'analyse technique porte sur la robustesse des appels HTTP, la gestion des erreurs, la sécurité des endpoints et la performance des opérations. Les interceptors Angular jouent un rôle central : l'interceptor d'authentification ajoute le token, l'interceptor d'erreur peut intercepter les réponses non autorisées et déclencher une déconnexion ou un message utilisateur.

\subsubsection*{Architecture des composants concernés}

L'architecture finale repose sur une communication REST entre Angular et ASP.NET Core. Angular fournit les pages et services de présentation, tandis que l'API centralise les règles métier et l'accès aux données. SQL Server conserve les utilisateurs, les configurations, les modèles et les données métier. Cette architecture facilite la maintenance, car chaque couche dispose de responsabilités identifiables.

\section{Réalisation}

\subsection{Validation des parcours utilisateurs}

Les parcours utilisateurs validés couvrent les trois rôles. Le super administrateur se connecte, accède à l'espace global, configure les familles et les vues de tables. L'administrateur gère les modèles, modifie les contenus et associe les chartes graphiques. L'utilisateur final consulte les familles disponibles, sélectionne un modèle, choisit un bénéficiaire et prévisualise un document.

\screenshotplaceholder{sprint5_demo.png}{Parcours de démonstration finale}{Démonstration d'un parcours complet de génération documentaire}

La validation a également porté sur les redirections automatiques. Un utilisateur déjà connecté ne doit pas revenir sur la page de connexion. Un utilisateur non autorisé ne doit pas accéder à un espace qui ne correspond pas à son rôle.

\subsection{Gestion des erreurs}

La gestion des erreurs est indispensable pour éviter les blocages silencieux. Côté Angular, les interceptors permettent de centraliser une partie du traitement. Côté backend, les contrôleurs retournent des réponses structurées lorsque des paramètres obligatoires sont absents ou lorsqu'une ressource n'existe pas.

\begin{lstlisting}[caption={Principe d'un interceptor d'erreur Angular}]
export const errorInterceptor: HttpInterceptorFn = (req, next) => {
  const authService = inject(AuthService);
  return next(req).pipe(
    catchError((error) => {
      if (error.status === 401) {
        authService.logout();
      }
      return throwError(() => error);
    })
  );
};
\end{lstlisting}

Ce type de mécanisme améliore l'expérience utilisateur et évite qu'une session expirée produise des erreurs répétées dans les pages métier.

\subsection{Optimisation et renforcement de la sécurité}

Le renforcement de la sécurité s'est appuyé sur plusieurs vérifications. Les endpoints critiques doivent être protégés par JWT. Les données envoyées par le client doivent être validées. Les requêtes SQL dynamiques doivent limiter les champs manipulés et utiliser des paramètres. Les mots de passe ne doivent jamais être exposés dans les réponses HTTP.

Les optimisations portent également sur le chargement côté frontend. Les routes utilisent le chargement paresseux des composants, ce qui évite de charger immédiatement toutes les pages. Les services d'état évitent des appels API répétés en conservant les ressources déjà chargées.

\subsection{Documentation technique et utilisateur}

La documentation technique présente l'architecture générale, les endpoints API, les entités principales, les choix technologiques et les diagrammes UML. La documentation utilisateur décrit les principales actions : connexion, navigation selon le rôle, création d'une famille, création d'un modèle, configuration d'une charte graphique, configuration d'une vue de table et prévisualisation d'un document.

Les captures d'écran doivent être insérées dans le dossier \texttt{images/}. Les diagrammes Draw.io doivent être conservés dans le dossier \texttt{diagrammes/}, puis exportés en PNG pour être inclus dans le rapport final.

\subsection{Préparation de la démonstration PFE}

La démonstration PFE doit suivre un scénario clair et progressif. Elle commence par la connexion d'un super administrateur, se poursuit par la présentation des modules d'administration, puis montre l'édition d'un modèle et se termine par la prévisualisation d'un document généré à partir d'un bénéficiaire.

Le scénario recommandé est le suivant :

\begin{enumerate}
    \item connexion avec un compte de démonstration ;
    \item présentation rapide du tableau de bord ou de l'espace correspondant au rôle ;
    \item création ou consultation d'une famille de documents ;
    \item consultation d'un modèle associé à cette famille ;
    \item modification mineure du contenu dans l'éditeur Tiptap ;
    \item association ou visualisation d'une charte graphique ;
    \item sélection d'un bénéficiaire ;
    \item génération de la prévisualisation ;
    \item explication de l'architecture backend/frontend et des contrôles de sécurité.
\end{enumerate}

\screenshotplaceholder{sprint5_tests.png}{Résultats de tests et validation technique}{Résultat des tests, build ou validation finale}

\section{Tests et validation}

\begin{longtable}{p{0.25\textwidth} p{0.43\textwidth} p{0.22\textwidth}}
\caption{Scénarios de test du sprint 5}\\
\toprule
\textbf{Scénario} & \textbf{Description} & \textbf{Résultat attendu}\\
\midrule
\endfirsthead
\toprule
\textbf{Scénario} & \textbf{Description} & \textbf{Résultat attendu}\\
\midrule
\endhead
Parcours super administrateur & Connexion, accès espace global, gestion famille et vue table. & Parcours terminé sans erreur bloquante.\\
\midrule
Parcours administrateur & Connexion, modification modèle, association charte. & Modèle sauvegardé et rendu mis à jour.\\
\midrule
Parcours utilisateur final & Connexion, choix famille, choix bénéficiaire, prévisualisation. & Document prévisualisé correctement.\\
\midrule
Exploitation module métier & Accès à l'espace données, sélection d'un module et bascule entre les tables. & Affichage cohérent des tables regroupées.\\
\midrule
Expiration session & Utilisation d'un token expiré ou supprimé. & Redirection vers la connexion.\\
\midrule
CRUD complet & Création, modification et suppression de ressources métier. & Données persistées et état frontend cohérent.\\
\midrule
Build frontend & Compilation Angular en mode production. & Build terminé sans erreur critique.\\
\midrule
Build backend & Compilation du projet ASP.NET Core. & Build terminé sans erreur critique.\\
\bottomrule
\end{longtable}

Les tests de sécurité ont vérifié le refus des appels anonymes sur les endpoints protégés, l'absence du mot de passe haché dans les réponses et la redirection des utilisateurs vers leur espace autorisé. Les tests d'intégration ont validé les flux entre Angular, l'API et SQL Server. Les tests fonctionnels ont confirmé que les principaux scénarios de démonstration étaient réalisables.

\section{Conclusion}

Le sprint 5 a permis de consolider l'ensemble du projet et de préparer sa livraison. Les modules développés dans les sprints précédents ont été intégrés dans des parcours cohérents. Les tests ont permis d'identifier les points critiques, de vérifier les contrats API et de valider les scénarios nécessaires à la démonstration.

À l'issue de ce sprint, la plateforme dispose d'une base fonctionnelle complète : authentification, rôles, édition documentaire, administration métier, configuration des vues de tables, gestion des bénéficiaires et prévisualisation dynamique. Le projet est ainsi prêt à être présenté comme une solution professionnelle de gestion documentaire dynamique.
